\chapter{Probability Spaces}\label{ch:probability-spaces}

\chapterinfo{STAT2003, STAT1301}{\textit{none}}

% Sources: STAT2003 handwritten notes (pp. 1-7, 84); STAT2003 course notes, Chapter 1;
% STAT1301 course notes (2026), Sections 1.3-1.7.

The basic notion in probability is that of a random experiment. This chapter builds the mathematical model of a random experiment, which has three ingredients: a sample space, a collection of events and a probability. It then covers counting, which is how most probabilities on finite sample spaces are calculated, and finally conditional probability and independence.

% ==========================================================================
\section{Random experiments}\label{sec:random-experiments}
% ==========================================================================

A \emph{random experiment} is an experiment whose outcome cannot be determined in advance, but which is nevertheless subject to analysis. Examples of random experiments are:
\begin{enumerate}
  \item tossing a die;
  \item measuring the amount of rainfall in Brisbane in January;
  \item counting the number of calls arriving at a telephone exchange during a fixed time period;
  \item selecting a random sample of fifty people and observing the number of left-handers;
  \item choosing ten people at random and measuring their heights;
  \item conducting an observational study on post-heart-failure patients, which results in a set of measurements for each patient.
\end{enumerate}
The word \emph{experiment} often means a planned scientific study, but here it is used in a broader sense, for any situation whose outcome is treated as random.

The goal of probability is to understand the behaviour of random experiments by analysing the corresponding mathematical models. Given a model, we can calculate quantities of interest such as probabilities and expectations. Mathematical models of random experiments are also the basis of statistics, where an appropriate model is fitted to observed data.

\begin{example}[Coin tossing]\label{ex:coin-tossing-simulation}
The most fundamental random experiment is tossing a coin a number of times. Much of probability theory can be based on this simple experiment. To see how it behaves, we can carry it out on a computer. The following R code simulates 100 tosses of a fair coin (Heads and Tails equally likely), plots the outcomes as a bar chart, and then plots the average number of Heads against the number of tosses.
\begin{rcode}
x <- sample(c(0, 1), size = 100, replace = TRUE)  # simulate coin tosses
barplot(x)                    # bar chart of the outcomes (1 = Heads)
y <- cumsum(x) / 1:100        # running average number of Heads
plot(y, type = "l")           # line graph of the running average
\end{rcode}
Repeating this a few times shows that the average number of Heads seems to converge to $\tfrac12$, but with a lot of random fluctuation. Similar results hold for a biased coin with probability $p$ of Heads.

Typical \emph{probability} questions about this experiment are: what is the probability of $x$ Heads in 100 tosses? How many Heads would you expect? What is the probability of waiting more than 4 tosses for the first Head? A \emph{statistical} analysis instead starts from observed data, such as the outcomes of 100 tosses, with $p$ unknown. Typical questions are: is the coin fair? How can $p$ best be estimated from the data? How accurate is that estimate?
\end{example}

Some further examples of random experiments:
\begin{itemize}
  \item \textbf{Control charts.} In manufacturing, the average output of a process is measured every hour, for example the average weight of 10 bags of sugar, to assess whether the process is still ``in control''. If the average goes above an upper control limit or below a lower control limit, an alarm is raised and the machine is adjusted. The question is how to set the control limits, since the process naturally fluctuates around its target.
  \item \textbf{Machine lifetimes.} 1000 identical components are monitored for failure, up to 50\,000 hours. The outcome is summarised in a cumulative lifetime table giving, for each time $t$, the proportion $\widehat{F}(t)$ of components that have failed by time $t$. One question is how $\widehat{F}$ can be modelled by a continuous function $F$, the lifetime distribution of a typical component.
  \item \textbf{An aeroplane.} A four-engine aeroplane can fly on just one engine on each wing, and all engines are unreliable. Number the engines 1, 2 (left wing) and 3, 4 (right wing), and observe which engines work during a specified period. There are $2^4 = 16$ possible outcomes. Which outcomes lead to system failure? If the probability of failure is known for each engine, what is the probability of failure for the whole system?
\end{itemize}

We describe random experiments by mathematical models with three building blocks: a sample space, a set of events, and a probability.

% ==========================================================================
\section{Sample spaces and events}\label{sec:sample-spaces-events}
% ==========================================================================

Although we cannot predict the outcome of a random experiment with certainty, we can usually specify the set of possible outcomes.

\begin{definition}[Sample space]\label{def:sample-space}
The \emph{sample space} $\Omega$ of a random experiment is the set of all possible outcomes of the experiment.
\end{definition}

\begin{example}[Sample spaces]\label{ex:sample-spaces}
\leavevmode
\begin{enumerate}
  \item Cast two dice consecutively: $\Omega = \set{(1,1), (1,2), \ldots, (1,6), (2,1), \ldots, (6,6)}$.
  \item The lifetime of a machine, in days: $\Omega = \R_+$, the positive real numbers.
  \item The number of calls arriving at an exchange during a specified time interval: $\Omega = \set{0, 1, \ldots} = \Z_+$.
  \item The heights of 10 selected people: $\Omega = \set{(x_1, \ldots, x_{10}) : x_i \ge 0, \ i = 1, \ldots, 10} = \R_+^{10}$. Here $x_1$ is the height of the first selected person, $x_2$ that of the second, and so on.
\end{enumerate}
For modelling purposes, it is often easier to take the sample space larger than necessary. The actual lifetime of a machine would certainly not span the whole positive real line, and the heights of 10 people would not exceed 3 metres.
\end{example}

Often we are not interested in a single outcome but in whether or not one of a group of outcomes occurs.

\begin{definition}[Event]\label{def:event}
An \emph{event} is a subset of the sample space. Events are denoted by capital letters $A, B, C, \ldots$. An event $A$ \emph{occurs} if the outcome of the experiment is one of the elements of $A$.
\end{definition}

\begin{example}[Events]\label{ex:events}
\leavevmode
\begin{enumerate}
  \item The event that the sum of two dice is 10 or more: $A = \set{(4,6), (5,5), (5,6), (6,4), (6,5), (6,6)}$.
  \item The event that a machine lives less than 1000 days: $A = [0, 1000)$.
  \item The event that, out of fifty selected people, five are left-handed: $A = \set{5}$.
  \item A coin is tossed 3 times and every head and tail is recorded. The sample space is
  \[
    \Omega = \set{HHH, HHT, HTH, HTT, THH, THT, TTH, TTT},
  \]
  where, for example, $HTH$ means heads, then tails, then heads. An alternative sample space is the set $\set{0,1}^3$ of binary vectors of length 3, so that $HTH$ corresponds to $(1,0,1)$. The event that the third toss is heads is $A = \set{HHH, HTH, THH, TTH}$.
\end{enumerate}
\end{example}

Since events are sets, we can apply the usual set operations to them:
\begin{enumerate}
  \item $A \cup B$ (\emph{$A$ union $B$}) is the event that $A$ or $B$ or both occur;
  \item $A \cap B$ (\emph{$A$ intersection $B$}) is the event that $A$ and $B$ both occur;
  \item $A^c$ (\emph{$A$ complement}) is the event that $A$ does not occur;
  \item if $A \subseteq B$ ($A$ is a subset of $B$), then $A$ is said to \emph{imply} $B$.
\end{enumerate}
Two events $A$ and $B$ with no outcomes in common, $A \cap B = \emptyset$, are \emph{disjoint}. Events are often pictured in a Venn diagram.

\begin{example}[Casting two dice]\label{ex:two-dice-intersection}
Cast two dice consecutively, with $\Omega$ as in \cref{ex:sample-spaces}. Let $A = \set{(6,1), \ldots, (6,6)}$ be the event that the first die is 6, and $B = \set{(1,6), \ldots, (6,6)}$ the event that the second die is 6. Then $A \cap B = \set{(6,6)}$ is the event that both dice are 6.
\end{example}

\begin{theorem}[Set properties and laws]\label{thm:set-laws}
Let $A$, $B$ and $C$ be sets. Then
\begin{align*}
  &\text{Commutative laws:} & A \cup B &= B \cup A, & A \cap B &= B \cap A, \\
  &\text{Associative laws:} & A \cup (B \cup C) &= (A \cup B) \cup C, & A \cap (B \cap C) &= (A \cap B) \cap C, \\
  &\text{Distributive laws:} & A \cup (B \cap C) &= (A \cup B) \cap (A \cup C), & A \cap (B \cup C) &= (A \cap B) \cup (A \cap C), \\
  &\text{De Morgan's laws:} & (A \cap B)^c &= A^c \cup B^c, & (A \cup B)^c &= A^c \cap B^c.
\end{align*}
\end{theorem}

De Morgan's laws hold for any collection of events $\set{A_i}$:
\begin{equation}\label{eq:de-morgan}
  \Bigl(\bigcup_i A_i\Bigr)^c = \bigcap_i A_i^c, \qquad \Bigl(\bigcap_i A_i\Bigr)^c = \bigcup_i A_i^c.
\end{equation}
These are easily checked with Venn diagrams. If $A_i$ is the event that component $i$ of a system works, the left-hand side of the first law is the event that a parallel system (which works if at least one component works) is not working. The right-hand side is the event that all components are not working. Clearly these are the same event.

\begin{example}[System reliability]\label{ex:system-reliability-events}
Consider three systems, each made of 3 unreliable components. The \emph{series} system works if and only if all components work; the \emph{parallel} system works if and only if at least one component works; and the \emph{2-out-of-3} system works if and only if at least 2 of the 3 components work. Let $A_i$ be the event that component $i$ works, and $D_a$, $D_b$, $D_c$ the events that the series, parallel and 2-out-of-3 systems work. Then
\begin{align*}
  D_a &= A_1 \cap A_2 \cap A_3, \qquad D_b = A_1 \cup A_2 \cup A_3, \\
  D_c &= (A_1 \cap A_2 \cap A_3) \cup (A_1^c \cap A_2 \cap A_3) \cup (A_1 \cap A_2^c \cap A_3) \cup (A_1 \cap A_2 \cap A_3^c) \\
      &= (A_1 \cap A_2) \cup (A_1 \cap A_3) \cup (A_2 \cap A_3).
\end{align*}
\end{example}

% ==========================================================================
\section{Probability}\label{sec:probability}
% ==========================================================================

The third ingredient of the model is the probability of the events, which says how likely each event is to occur.

\begin{definition}[Probability space]\label{def:probability-space}
A \emph{probability space} is a triple $(\Omega, \mathcal{F}, \Prob)$, where
\begin{itemize}
  \item $\Omega$ is the sample space;
  \item $\mathcal{F} \subseteq 2^\Omega$ is the collection of events, where $2^\Omega$ is the set of all subsets of $\Omega$;
  \item $\Prob \colon \mathcal{F} \to [0, 1]$ is the probability (also called the probability law or probability measure).
\end{itemize}
\end{definition}

\begin{remark}
If $\Omega$ is infinite, $\mathcal{F}$ is infinite too. In that case $\mathcal{F}$ is taken to be a $\sigma$-algebra.
\end{remark}

\begin{definition}[Probability measure]\label{def:probability-measure}
A \emph{probability} $\Prob$ is a rule (function) that assigns a number to each event and satisfies the following axioms (Kolmogorov, 1933):
\begin{enumerate}[label=\textbf{Axiom \arabic*:}, leftmargin=*, widest=3]
  \item $\Prob(A) \ge 0$ for every event $A$;
  \item $\Prob(\Omega) = 1$;
  \item for any sequence $A_1, A_2, \ldots$ of disjoint events,
  \begin{equation}\label{eq:sum-rule}
    \Prob\Bigl(\bigcup_i A_i\Bigr) = \sum_i \Prob(A_i).
  \end{equation}
\end{enumerate}
\end{definition}

Axiom 2 says that the ``certain'' event $\Omega$ has probability 1. Property \eqref{eq:sum-rule} is the crucial property of a probability, and is called the \emph{sum rule}: if an event can happen in a number of different ways that cannot happen at the same time, its probability is the sum of the probabilities of those ways. A probability behaves exactly like an area measure, where the area of a union of non-overlapping regions is the sum of their areas. Instead of measuring areas, $\Prob$ measures probabilities.

\begin{example}[An unfair coin]\label{ex:unfair-coin}
A coin with heads less likely than tails: $\Omega = \set{H, T}$, $\mathcal{F} = 2^\Omega = \set{\emptyset, \set{H}, \set{T}, \set{H, T}}$, and
\[
  \Prob(\set{H}) = 0.45, \qquad \Prob(\set{T}) = 0.55, \qquad \Prob(\emptyset) = 0, \qquad \Prob(\Omega) = 1.
\]
\end{example}

\begin{theorem}[Properties of a probability measure]\label{thm:probability-properties}
Let $A$ and $B$ be events. Then
\begin{enumerate}
  \item $\Prob(\emptyset) = 0$;
  \item $A \subseteq B \implies \Prob(A) \le \Prob(B)$;
  \item $\Prob(A) \le 1$;
  \item $\Prob(A^c) = 1 - \Prob(A)$ (complement rule);
  \item $\Prob(A \cup B) = \Prob(A) + \Prob(B) - \Prob(A \cap B)$ (addition rule).
\end{enumerate}
\end{theorem}

\begin{proof}
1. $\Omega = \Omega \cup \emptyset$, with $\Omega$ and $\emptyset$ disjoint, so by the sum rule $\Prob(\Omega) = \Prob(\Omega) + \Prob(\emptyset)$. By Axiom 2, $1 = 1 + \Prob(\emptyset)$, so $\Prob(\emptyset) = 0$.

2. If $A \subseteq B$, then $B = A \cup (B \setminus A)$, where $A$ and $B \setminus A$ are disjoint. By the sum rule, $\Prob(B) = \Prob(A) + \Prob(B \setminus A)$, which is at least $\Prob(A)$ by Axiom 1.

3. $A \subseteq \Omega$, so $\Prob(A) \le \Prob(\Omega) = 1$ by property 2 and Axiom 2.

4. $\Omega = A \cup A^c$, where $A$ and $A^c$ are disjoint. By the sum rule and Axiom 2, $1 = \Prob(\Omega) = \Prob(A) + \Prob(A^c)$.

5. $A \cup B$ is the disjoint union of $A$ and $B \cap A^c$, so $\Prob(A \cup B) = \Prob(A) + \Prob(B \cap A^c)$. Also, $B$ is the disjoint union of $A \cap B$ and $B \cap A^c$, so $\Prob(B) = \Prob(A \cap B) + \Prob(B \cap A^c)$. Combining the two equations gives the result.
\end{proof}

The addition rule extends to $n$ events.

\begin{theorem}[Inclusion--exclusion principle]\label{thm:inclusion-exclusion}
For events $E_1, \ldots, E_n$,
\[
  \Prob\Bigl(\bigcup_{i=1}^n E_i\Bigr) = \sum_i \Prob(E_i) - \sum_{i<j} \Prob(E_i \cap E_j) + \sum_{i<j<k} \Prob(E_i \cap E_j \cap E_k) - \cdots + (-1)^{n+1} \Prob(E_1 \cap \cdots \cap E_n).
\]
In particular, $\Prob\bigl(\bigcup_i E_i\bigr) \le \sum_i \Prob(E_i)$.
\end{theorem}

The model for a random experiment is now complete. It is up to the modeller to choose the sample space $\Omega$ and probability $\Prob$ that best describe the actual experiment.

\begin{example}[Fair die]\label{ex:fair-die}
Throw a fair die. Then $\Omega = \set{1, 2, \ldots, 6}$, and common sense says that
\[
  \Prob(A) = \frac{\abs{A}}{6}, \qquad A \subseteq \Omega,
\]
where $\abs{A}$ is the number of elements of $A$. For example, the probability of an even number is $\Prob(\set{2,4,6}) = \tfrac36 = \tfrac12$.
\end{example}

\subsection{Discrete and continuous sample spaces}\label{sec:discrete-continuous-sample-spaces}

A sample space is \emph{discrete} if it is countable, that is, $\Omega = \set{a_1, \ldots, a_n}$ or $\Omega = \set{a_1, a_2, \ldots}$. The easiest way to specify a probability on a discrete sample space is to give the probability $p_i$ of each elementary event $\set{a_i}$, and then set
\[
  \Prob(A) = \sum_{i : a_i \in A} p_i \qquad \text{for all } A \subseteq \Omega.
\]
In many applications all elementary events are equally likely, so each has probability $1/\abs{\Omega}$. In \cref{ex:fair-die}, each has probability $\tfrac16$. This principle is important enough to state as a theorem.

\begin{theorem}[Equilikely principle]\label{thm:equilikely-principle}
If $\Omega$ has a finite number of outcomes, all equally likely, then the probability of each event $A$ is
\[
  \Prob(A) = \frac{\abs{A}}{\abs{\Omega}} = \frac{\text{number of elements of } A}{\text{number of elements of } \Omega}.
\]
\end{theorem}

For such sample spaces, calculating probabilities reduces to counting (\cref{sec:counting}).

A sample space that is not countable, such as $\Omega = \R_+$, is \emph{continuous}.

\begin{example}[Drawing a point from an interval]\label{ex:point-in-interval}
Draw a point at random from $[0,1]$, each point being equally likely. The sample space is $\Omega = [0, 1]$, which is continuous. We cannot define $\Prob$ through the elementary events $\set{x}$, $x \in [0,1]$, because each of them must have probability $0$. Instead, define
\[
  \Prob([a, b]) = b - a, \qquad 0 \le a \le b \le 1.
\]
This completely specifies $\Prob$. In particular, the probability that the point falls in any (sufficiently nice) set $A$ is the length of $A$.
\end{example}

% ==========================================================================
\section{Counting}\label{sec:counting}
% ==========================================================================

Counting is not always easy. Consider these examples:
\begin{enumerate}
  \item A multiple-choice form has 20 questions, each with 3 choices. In how many ways can the form be completed?
  \item In a horse race with 8 horses, how many ways are there to bet on the placings (1st, 2nd, 3rd)?
  \item Jessica has a collection of 20 CDs and wants to take 3 of them to work. How many possibilities does she have?
  \item How many different throws are possible with 3 dice?
\end{enumerate}
Solving counting problems comfortably takes a lot of practice. Fortunately, many of them fit the simple framework of drawing balls from an urn.

Consider an urn with $n$ different balls, numbered $1, \ldots, n$, from which $k$ balls are drawn. The balls can be drawn one by one, noting the order in which they are drawn, or the order can be ignored (as if all $k$ balls were drawn in one grab). And once a ball is drawn, it can either be put back in the urn after its number is recorded, or left out: drawing \emph{with} or \emph{without replacement}. This gives four possible experiments:
\begin{enumerate}
  \item ordered, with replacement;
  \item ordered, without replacement;
  \item unordered, without replacement;
  \item unordered, with replacement.
\end{enumerate}
The art is to recognise a seemingly unrelated counting problem as one of these four urn problems. Example 1 above is drawing 20 balls from an urn of 3, ordered, with replacement. Example 2 is drawing 3 balls from an urn of 8, ordered, without replacement. Example 3 is drawing 3 balls from an urn of 20, unordered, without replacement. Example 4 is drawing 3 balls from an urn of 6, unordered, with replacement.

Ordered arrangements are written as vectors with round brackets, for example $(1,2,3) \ne (3,2,1)$. Unordered arrangements are written as sets with curly brackets, for example $\set{1,2,3} = \set{3,2,1}$. In each case below, we first count for $n = 4$ and $k = 3$ and then state the general result.

\paragraph{Ordered, with replacement.} After each ball is drawn, its number is noted and the ball is put back. Possible outcomes for $n=4$, $k=3$ include $(1,1,1)$, $(4,1,2)$, $(2,3,2)$ and $(4,2,1)$. There are three positions $(\cdot, \cdot, \cdot)$ to fill, and each can hold 1, 2, 3 or 4, so there are $4 \times 4 \times 4 = 4^3 = 64$ possibilities. This can be pictured as a tree diagram, with one level per position.

\begin{theorem}[Ordered arrangements with replacement]\label{thm:count-ordered-with-replacement}
The number of ordered arrangements of $k$ numbers chosen from $\set{1, \ldots, n}$ with replacement (repetition allowed) is $n^k$.
\end{theorem}

\paragraph{Ordered, without replacement.} Now the numbers in the three positions must be different, for example $(1,4,2)$ or $(3,2,1)$. There are 4 possibilities for the first position, then 3 for the second, then 2 for the third, so there are $4 \times 3 \times 2 = 24$ arrangements. This is the tree for the previous case with all duplicate branches removed. An ordered arrangement without repetition is called a \emph{permutation} of size $k$ from $\set{1, \ldots, n}$. A permutation of size $n$ is simply called a permutation of $\set{1, \ldots, n}$.

\begin{theorem}[Ordered arrangements without replacement]\label{thm:count-ordered-without-replacement}
The number of permutations of size $k$ from $\set{1, \ldots, n}$ is
\[
  {}^nP_k = n(n-1)\cdots(n-k+1) = \frac{n!}{(n-k)!}.
\]
In particular, there are $n! = n(n-1)(n-2)\cdots 1$ ($n$ factorial) ordered arrangements of $n$ items.
\end{theorem}

\paragraph{Unordered, without replacement.} Now the order is not noted either, so possible outcomes for $n=4$, $k=3$ are $\set{1,2,4}$, $\set{1,2,3}$ and so on. If we noted the order, there would be ${}^nP_k$ outcomes, among them $(1,2,4)$, $(1,4,2)$, $(2,1,4)$, $(2,4,1)$, $(4,1,2)$ and $(4,2,1)$. These 6 permutations all correspond to the single unordered arrangement $\set{1,2,4}$. Unordered arrangements without repetition are called \emph{combinations} of size $k$ from $\set{1, \ldots, n}$. To count them, divide ${}^nP_k$ by the number of permutations of $k$ items, $k!$. For $n=4$, $k=3$ there are $24/6 = 4$ combinations.

\begin{theorem}[Unordered arrangements without replacement]\label{thm:count-unordered-without-replacement}
The number of combinations of size $k$ from $\set{1, \ldots, n}$ is
\[
  {}^nC_k = \binom{n}{k} = \frac{{}^nP_k}{k!} = \frac{n!}{(n-k)!\,k!}.
\]
\end{theorem}

\paragraph{Unordered, with replacement.} For $n=4$, $k=3$, possible outcomes are $\set{3,3,4}$, $\set{1,2,4}$, $\set{2,2,2}$ and so on. The trick is to represent each outcome by a vector $(x_1, \ldots, x_n)$ recording how many times each number occurs. For example, $\set{3,3,4}$ corresponds to $(0,0,2,1)$ and $\set{1,2,4}$ to $(1,1,0,1)$. So we must count the vectors with non-negative integer components summing to $k$. Equivalently, place $k = 3$ balls into $n = 4$ urns numbered $1, \ldots, 4$: the vector $(0,0,2,1)$ means urn 3 has 2 balls and urn 4 has 1. One way to do this is to distribute $n - 1 = 3$ ``separators'' and $k = 3$ balls over $n - 1 + k = 6$ positions. The number of ways is the number of ways to choose the $k$ positions for the balls out of $n - 1 + k$.

\begin{theorem}[Unordered arrangements with replacement]\label{thm:count-unordered-with-replacement}
The number of different multisets $\set{x_1, \ldots, x_k}$ with $x_i \in \set{1, \ldots, n}$ is
\[
  \binom{n + k - 1}{k}.
\]
\end{theorem}

\begin{remark}[Balls into bins]\label{rem:balls-into-bins}
The urn framework is equivalent to putting $k$ balls into $n$ bins:
\begin{center}
\begin{tabular}{@{}lll@{}}
\toprule
Urn: draw $k$ of $n$ balls & Bins: put $k$ balls in $n$ bins & Number of ways \\
\midrule
ordered, with replacement & balls differ, no limit per bin & $n^k$ \\
ordered, without replacement & balls differ, at most 1 per bin & ${}^nP_k$ \\
unordered, without replacement & balls identical, at most 1 per bin & $\binom{n}{k}$ \\
unordered, with replacement & balls identical, no limit per bin & $\binom{n+k-1}{k}$ \\
\bottomrule
\end{tabular}
\end{center}
\end{remark}

\begin{example}[Tossing a pair of dice]\label{ex:counting-pair-of-dice}
Count the outcomes of tossing two dice ($n = 6$, $k = 2$) under the four schemes.
\end{example}

\begin{solution}
\begin{center}
\begin{tabular}{@{}lcc@{}}
\toprule
 & repeats allowed & repeats not allowed \\
\midrule
dice of different colours (ordered) & $6^2 = 36$ & ${}^6P_2 = 30$ \\
dice of the same colour (unordered) & $\binom{7}{2} = 21$ & $\binom{6}{2} = 15$ \\
\bottomrule
\end{tabular}
\end{center}
\end{solution}

Returning to the four problems at the start of the section:
\begin{enumerate}
  \item The form can be completed in $3^{20} = 3\,486\,784\,401$ ways.
  \item The number of placings is ${}^8P_3 = 336$.
  \item The number of possible combinations of CDs is $\binom{20}{3} = 1140$.
  \item The number of different throws with three dice is $\binom{8}{3} = 56$.
\end{enumerate}

Not every problem reduces directly to one of the four urn problems, but the counting principles are the same. Some more examples:
\begin{enumerate}
  \item The numbers $1, \ldots, 5$ can be arranged (as in 13524, 25134, \ldots) in $5! = 120$ ways.
  \item The arrangements of $1, 2, \ldots, 7$ in which the first 3 numbers are 1, 2, 3 (in any order) and the last 4 are 4, 5, 6, 7 (in any order) number $3! \times 4!$.
  \item The word ``arrange'' can be rearranged (as in ``aarrnge'', ``arrngea'', \ldots) in $7!/4 = 1260$ ways. Label the letters as 7 balls: balls 1 and 2 are `a', 3 and 4 are `r', 5 is `n', 6 is `g', 7 is `e'. There are $7!$ permutations of the balls, but swapping the two `a's or the two `r's gives the same word, so divide by $2! \times 2! = 4$.
  \item An urn has 1000 balls labelled $000, 001, \ldots, 999$. The number of balls whose digits are in strictly ascending order (such as 047 and 489, but not 033 or 321) is $120$. There are $10 \times 9 \times 8 = 720$ labels with three different digits. Each set of three digits can be arranged in $3! = 6$ ways, only one of which is ascending, so the answer is $720/6 = 120$.
  \item In a group of 20 people, each with a different birthday, the number of possible arrangements of the birthdays (with 365 days in a year) is ${}^{365}P_{20}$.
\end{enumerate}

Once we can count, the equilikely principle (\cref{thm:equilikely-principle}) gives probabilities.

\begin{example}[Probabilities by counting]\label{ex:probabilities-by-counting}
\leavevmode
\begin{enumerate}
  \item \emph{What is the probability that, in a group of 40 people, all have different birthdays?} Choosing the birthdays is like drawing 40 balls with replacement from an urn with balls $1, \ldots, 365$. So $\Omega$ consists of vectors of length 40 with components in $\set{1, \ldots, 365}$. There are $\abs{\Omega} = 365^{40}$ of them, all equally likely. Let $A$ be the event that all 40 people have different birthdays. Then $\abs{A} = {}^{365}P_{40} = 365!/325!$, so $\Prob(A) = \abs{A}/\abs{\Omega} \approx 0.109$. Not very big!
  \item \emph{What is the probability of exactly 5 Heads and 5 Tails in 10 tosses of a fair coin?} $\Omega$ consists of the $2^{10}$ binary vectors of length 10 (1 for Heads), all equally likely. The event $A$ of exactly 5 Heads contains as many vectors as there are ways to choose the positions of the five 1s out of 10, namely $\binom{10}{5}$. So $\Prob(A) = \binom{10}{5}/2^{10} = 252/1024 \approx 0.25$.
  \item \emph{13 cards are drawn at random from a full deck. What is the probability of 4 Hearts and 3 Diamonds?} Number the cards 1 to 52, with 1--13 Hearts, 14--26 Diamonds, and so on. $\Omega$ consists of the unordered sets of 13 cards without repetition. There are $\abs{\Omega} = \binom{52}{13}$ of them, all equally likely. To form the event $A$, choose 4 Hearts out of 13, 3 Diamonds out of 13, and the other 6 cards out of the 26 Spades and Clubs: $\abs{A} = \binom{13}{4}\binom{13}{3}\binom{26}{6}$. So $\Prob(A) = \abs{A}/\abs{\Omega} \approx 0.074$.
\end{enumerate}
\end{example}

% ==========================================================================
\section{Conditional probability}\label{sec:conditional-probability}
% ==========================================================================

How do probabilities change when we know that some event $B \subseteq \Omega$ has occurred? Then the outcome lies in $B$, so $A$ occurs if and only if $A \cap B$ occurs. The relative chance of $A$ occurring is therefore $\Prob(A \cap B)/\Prob(B)$.

\begin{definition}[Conditional probability]\label{def:conditional-probability}
Let $A, B \in \mathcal{F}$ with $\Prob(B) \ne 0$. The \emph{conditional probability of $A$ given $B$} is
\begin{equation}\label{eq:conditional-probability}
  \Prob(A \given B) = \frac{\Prob(A \cap B)}{\Prob(B)}.
\end{equation}
\end{definition}

\begin{proposition}\label{prop:conditional-probability-is-probability}
For fixed $B$ with $\Prob(B) \ne 0$, $\Prob(\cdot \given B)$ satisfies the probability axioms.
\end{proposition}

\begin{proof}
$\Prob(A \given B) = \Prob(A \cap B)/\Prob(B) \ge 0$, since both numerator and denominator are non-negative. $\Prob(\Omega \given B) = \Prob(B)/\Prob(B) = 1$. If $A_1, A_2, \ldots$ are disjoint, then so are $A_1 \cap B, A_2 \cap B, \ldots$, so by the sum rule $\Prob\bigl(\bigcup_i A_i \given B\bigr) = \sum_i \Prob(A_i \cap B)/\Prob(B) = \sum_i \Prob(A_i \given B)$.
\end{proof}

\begin{example}[Throwing two dice]\label{ex:two-dice-conditional}
Throw two fair dice. Given that the sum of the eyes is 10, what is the probability that one 6 is cast?
\end{example}

\begin{solution}
Let $B = \set{(4,6), (5,5), (6,4)}$ be the event that the sum is 10, and let
\[
  A = \set{(1,6), \ldots, (5,6), (6,1), \ldots, (6,5)}
\]
be the event that one 6 is cast. Then $A \cap B = \set{(4,6), (6,4)}$. All elementary events are equally likely, so
\[
  \Prob(A \given B) = \frac{2/36}{3/36} = \frac23.
\]
\end{solution}

\begin{example}[Monty Hall problem]\label{ex:monty-hall}
A quiz contestant chooses one of three curtains, A, B or C, behind one of which a prize is hidden. Suppose, without loss of generality, that the contestant chooses A. The quiz master, Monty Hall, then always opens one of the other curtains. If the prize is behind B, he opens C; if it is behind C, he opens B; and if it is behind A, he opens B or C with equal probability (by tossing a coin the contestant cannot see). Suppose, again without loss of generality, that Monty opens B. The contestant may now switch to curtain C. Should they stay with A or switch to C?
\end{example}

\begin{solution}
The sample space has 4 outcomes: $Ac$ (the prize is behind A and Monty opens C), $Ab$, $Bc$ and $Cb$. The events that the prize is behind A, B and C are $A = \set{Ac, Ab}$, $B = \set{Bc}$ and $C = \set{Cb}$. These are equally likely, and $Ab$ and $Ac$ are equally likely, so
\[
  \Prob(\set{Ab}) = \Prob(\set{Ac}) = \tfrac16, \qquad \Prob(\set{Bc}) = \Prob(\set{Cb}) = \tfrac13.
\]
Monty opening curtain B means the event $\set{Ab, Cb}$ has occurred. Hence
\[
  \Prob(A \given \text{B opened}) = \frac{\Prob(\set{Ac, Ab} \cap \set{Ab, Cb})}{\Prob(\set{Ab, Cb})} = \frac{1/6}{1/6 + 1/3} = \frac13,
\]
\[
  \Prob(C \given \text{B opened}) = \frac{\Prob(\set{Cb})}{\Prob(\set{Ab, Cb})} = \frac{1/3}{1/6 + 1/3} = \frac23.
\]
Opening a curtain gives no extra information about whether the prize is behind A. But $\Prob(B \given \text{B opened}) = 0$, and conditional probabilities must sum to 1, so the probability moves to C. Given that B is opened, the prize is twice as likely to be behind C as behind A, so the contestant should switch.
\end{solution}

\subsection{The product rule}\label{sec:product-rule}

By the definition of conditional probability,
\begin{equation}\label{eq:product-rule-two}
  \Prob(A \cap B) = \Prob(A)\,\Prob(B \given A).
\end{equation}
This generalises to $n$ events. We abbreviate $A_1 \cap A_2 \cap \cdots \cap A_n$ to $A_1 A_2 \cdots A_n$.

\begin{theorem}[Product rule]\label{thm:product-rule}
Let $A_1, \ldots, A_n$ be events with $\Prob(A_1 \cdots A_{n-1}) > 0$. Then
\begin{equation}\label{eq:product-rule}
  \Prob(A_1 \cdots A_n) = \Prob(A_1)\,\Prob(A_2 \given A_1)\,\Prob(A_3 \given A_1 A_2) \cdots \Prob(A_n \given A_1 \cdots A_{n-1}).
\end{equation}
\end{theorem}

\begin{proof}
We show the case of three events $A$, $B$, $C$; the case $n > 3$ is similar. Applying \eqref{eq:conditional-probability} to $\Prob(B \given A)$ and $\Prob(C \given A \cap B)$,
\[
  \Prob(A)\,\Prob(B \given A)\,\Prob(C \given A \cap B) = \Prob(A)\,\frac{\Prob(A \cap B)}{\Prob(A)}\,\frac{\Prob(A \cap B \cap C)}{\Prob(A \cap B)} = \Prob(A \cap B \cap C). \qedhere
\]
\end{proof}

\begin{example}[Drawing balls]\label{ex:drawing-balls-product-rule}
Three balls are drawn one after another, without replacement, from a bowl with 5 white and 5 black balls. What is the probability that all three are black?
\end{example}

\begin{solution}
Let $A_i$ be the event that the $i$th ball is black. By the product rule,
\[
  \Prob(A_1 A_2 A_3) = \Prob(A_1)\,\Prob(A_2 \given A_1)\,\Prob(A_3 \given A_1 A_2) = \frac{5}{10} \cdot \frac49 \cdot \frac38 \approx 0.083.
\]
This problem can also be solved by counting.
\end{solution}

\begin{example}[Birthday problem]\label{ex:birthday-product-rule}
By counting, the probability that 40 people all have different birthdays is
\begin{equation}\label{eq:birthday-40}
  \frac{365 \times 364 \times \cdots \times 326}{365 \times 365 \times \cdots \times 365} \approx 0.109
\end{equation}
(\cref{ex:probabilities-by-counting}). This also follows from the product rule. Let $A_i$ be the event that the first $i$ people have different birthdays. Then $A_1 \supseteq A_2 \supseteq \cdots$, so $A_n = A_1 \cap \cdots \cap A_n$, and
\[
  \Prob(A_{40}) = \Prob(A_1)\,\Prob(A_2 \given A_1)\,\Prob(A_3 \given A_2) \cdots \Prob(A_{40} \given A_{39}).
\]
Given that the first $k - 1$ people have different birthdays, there is no duplicate if and only if the $k$th person's birthday is one of the remaining $365 - (k-1)$ days, so $\Prob(A_k \given A_{k-1}) = (365 - k + 1)/365$. This gives \eqref{eq:birthday-40}. More generally, the probability that $n$ randomly chosen people have different birthdays is
\[
  \Prob(A_n) = \frac{365}{365} \times \frac{364}{365} \times \frac{363}{365} \times \cdots \times \frac{365 - n + 1}{365}, \qquad n \ge 1.
\]
This decreases rapidly to zero: for $n = 23$ it is already less than $\tfrac12$.
\end{example}

\subsection{The law of total probability and Bayes' rule}\label{sec:total-probability-bayes}

Events $B_1, B_2, \ldots, B_n$ form a \emph{partition} of $\Omega$ if they are disjoint and their union is $\Omega$.

\begin{theorem}[Law of total probability]\label{thm:total-probability}
Let $B_1, \ldots, B_n$ be a partition of $\Omega$ and $A$ an event. Then
\[
  \Prob(A) = \sum_{i=1}^n \Prob(A \given B_i)\,\Prob(B_i).
\]
\end{theorem}

\begin{proof}
$A$ is the disjoint union of $A \cap B_1, \ldots, A \cap B_n$, so by the sum rule $\Prob(A) = \sum_i \Prob(A \cap B_i)$. By the definition of conditional probability, $\Prob(A \cap B_i) = \Prob(A \given B_i)\,\Prob(B_i)$.
\end{proof}

\begin{theorem}[Bayes' rule]\label{thm:bayes-rule}
Let $B_1, \ldots, B_n$ be a partition of $\Omega$ and $A$ an event. Then
\[
  \Prob(B_j \given A) = \frac{\Prob(B_j \cap A)}{\Prob(A)} = \frac{\Prob(A \given B_j)\,\Prob(B_j)}{\sum_{i=1}^n \Prob(A \given B_i)\,\Prob(B_i)}.
\]
\end{theorem}

\begin{proof}
Combine the definition of conditional probability with the law of total probability.
\end{proof}

A simpler form of Bayes' rule is often appropriate. By the product rule, $\Prob(A \cap B) = \Prob(A)\,\Prob(B \given A) = \Prob(B)\,\Prob(A \given B)$, so
\[
  \Prob(B \given A) = \frac{\Prob(B)\,\Prob(A \given B)}{\Prob(A)}, \qquad \text{provided } \Prob(A) \ne 0.
\]

\begin{example}[Three factories]\label{ex:three-factories}
A company has three factories (1, 2 and 3) that produce the same chip, producing 15\%, 35\% and 50\% of the total production. The probability of a faulty chip at factories 1, 2 and 3 is 0.01, 0.05 and 0.02. A chip is selected at random from the total production and turns out to be faulty. What is the probability that it was produced in factory 1?
\end{example}

\begin{solution}
Let $B_i$ be the event that the chip was produced in factory $i$. The $B_i$ form a partition of $\Omega$. Let $A$ be the event that the chip is faulty. We are given $\Prob(B_1) = 0.15$, $\Prob(B_2) = 0.35$, $\Prob(B_3) = 0.5$ and $\Prob(A \given B_1) = 0.01$, $\Prob(A \given B_2) = 0.05$, $\Prob(A \given B_3) = 0.02$. By Bayes' rule,
\[
  \Prob(B_1 \given A) = \frac{0.15 \times 0.01}{0.15 \times 0.01 + 0.35 \times 0.05 + 0.5 \times 0.02} \approx 0.052.
\]
\end{solution}

% ==========================================================================
\section{Independence}\label{sec:independence}
% ==========================================================================

Independence is a very important concept in probability and statistics. Loosely speaking, it models a lack of information between events. $A$ and $B$ are independent if knowing that $A$ has occurred does not change the probability that $B$ occurs. Since $\Prob(A \given B) = \Prob(A \cap B)/\Prob(B)$, this is captured by the following definition, which also covers the case $\Prob(B) = 0$.

\begin{definition}[Independent events]\label{def:independent-events}
Events $A$ and $B$ are \emph{independent} if
\[
  \Prob(A \cap B) = \Prob(A)\,\Prob(B).
\]
When $\Prob(B) > 0$, this is equivalent to $\Prob(A \given B) = \Prob(A)$. More generally, the events $A_1, A_2, \ldots$ are \emph{(mutually) independent} if, for any $k$ and any choice of distinct indices $i_1, \ldots, i_k$,
\[
  \Prob(A_{i_1} \cap A_{i_2} \cap \cdots \cap A_{i_k}) = \Prob(A_{i_1})\,\Prob(A_{i_2}) \cdots \Prob(A_{i_k}).
\]
\end{definition}

\begin{remark}[Independence as a model assumption]\label{rem:independence-model-assumption}
In most cases, independence of events is a model assumption: we assume that there is a $\Prob$ under which certain events are independent.
\end{remark}

\begin{example}[Casting two dice, continued]\label{ex:two-dice-independent}
Cast two fair dice. Let $A$ be the event that the first die is 6 and $B$ the event that the second is 6. Naturally, $A$ and $B$ are independent: knowing that the first die is 6 gives no information about the second. Formally, $A = \set{(6,1), \ldots, (6,6)}$, $B = \set{(1,6), \ldots, (6,6)}$ and $A \cap B = \set{(6,6)}$, so
\[
  \Prob(A \given B) = \frac{1/36}{6/36} = \frac16 = \Prob(A).
\]
\end{example}

\begin{example}[Coin tosses and the binomial law]\label{ex:binomial-law}
Flip a coin $n$ times. The sample space is the set of binary $n$-tuples, $\Omega = \set{(0, \ldots, 0), \ldots, (1, \ldots, 1)}$, with 0 for Tails and 1 for Heads. For example, $(0, 1, 0, 1, \ldots)$ means Tails, Heads, Tails, Heads, and so on.

How should $\Prob$ be defined? Let $A_i$ be the event of Heads on the $i$th throw. $\Prob$ should make $A_1, \ldots, A_n$ independent, with $\Prob(A_i)$ the same for all $i$; call it $p \in [0, 1]$. These two requirements specify $\Prob$ completely. For example, the probability that the first $k$ throws are Heads and the last $n - k$ are Tails is
\[
  \Prob(\set{(1, \ldots, 1, 0, \ldots, 0)}) = \Prob(A_1) \cdots \Prob(A_k)\,\Prob(A_{k+1}^c) \cdots \Prob(A_n^c) = p^k (1-p)^{n-k}.
\]
Let $B_k$ be the event of $k$ Heads in total. $B_k$ is the union of the elementary events $\set{(x_1, \ldots, x_n)}$ with $x_1 + \cdots + x_n = k$. Each has probability $p^k(1-p)^{n-k}$, and there are $\binom{n}{k}$ of them, so
\[
  \Prob(B_k) = \binom{n}{k} p^k (1-p)^{n-k}, \qquad k = 0, 1, \ldots, n.
\]
This is the \emph{binomial distribution} (\cref{ch:common-distributions}).
\end{example}

\begin{example}[Geometric law]\label{ex:geometric-law}
Flip the same coin until Heads appears for the first time. Let $C_k$ be the event that the first Heads is on the $k$th toss, $k = 1, 2, \ldots$. With the events $A_i$ of \cref{ex:binomial-law},
\[
  C_k = A_1^c \cap A_2^c \cap \cdots \cap A_{k-1}^c \cap A_k.
\]
By the product rule and the independence of $A_1^c, \ldots, A_{k-1}^c, A_k$,
\[
  \Prob(C_k) = \Prob(A_1^c) \cdots \Prob(A_{k-1}^c)\,\Prob(A_k) = (1-p)^{k-1} p.
\]
This is the \emph{geometric law}.
\end{example}

\begin{example}[Cable matching]\label{ex:cable-matching}
$n$ cables, numbered $1, \ldots, n$ on one side, are connected at random to $n$ sockets numbered $1, \ldots, n$ on the other side. That is, they are connected by a random bijection, with all bijections equally likely. The system works if at least one cable $i$ is connected to socket $i$. What is the probability that the system works?
\end{example}

\begin{solution}
Let $E_i$ be the event that cable $i$ is connected properly. We need $\Prob(E_1 \cup \cdots \cup E_n)$, which we find by inclusion--exclusion (\cref{thm:inclusion-exclusion}). Here $\Prob(E_i) = \frac1n$, and by the product rule, for distinct indices,
\[
  \Prob(E_i \cap E_j) = \Prob(E_i)\,\Prob(E_j \given E_i) = \frac1n \cdot \frac{1}{n-1}, \qquad
  \Prob(E_i \cap E_j \cap E_k) = \frac1n \cdot \frac{1}{n-1} \cdot \frac{1}{n-2},
\]
and so on. There are $\binom{n}{k}$ intersections of $k$ distinct events, so
\[
  \sum_{i_1 < \cdots < i_k} \Prob(E_{i_1} \cap \cdots \cap E_{i_k}) = \binom{n}{k} \frac{1}{n(n-1)\cdots(n-k+1)} = \frac{1}{k!}.
\]
Therefore
\[
  \Prob\Bigl(\bigcup_{i=1}^n E_i\Bigr) = 1 - \frac{1}{2!} + \frac{1}{3!} - \cdots + \frac{(-1)^{n+1}}{n!} = \sum_{i=1}^n \frac{(-1)^{i+1}}{i!} \longrightarrow 1 - \frac1e \quad \text{as } n \to \infty.
\]
\end{solution}
